\documentclass[10pt,a4paper]{amsart}
\usepackage[margin=19mm]{geometry}
\usepackage{amsmath,amssymb,mathtools}
\usepackage[scr=rsfs,cal=euler]{mathalfa}
\usepackage{xfrac}
\usepackage[unicode,hidelinks,bookmarks=false]{hyperref}
\newcommand{\ie}{i.e.\ }
\newcommand{\cf}{cf.\ }
\newcommand{\resp}{respectively\ }
\newcommand{\Subsection}{\subsection}
\providecommand{\nobreakdash}{\nobreak\hbox{-}}
\setlength{\emergencystretch}{2em}
\multlinegap0pt
\usepackage{amssymb}
\usepackage{enumerate}
\newtheorem{theorem}{Theorem}
\newtheorem{lemma}{Lemma}
\def\supp{\hbox{supp\,}}
\title{Decay estimates for a class of dispersive equations with partial inverse-square potentials}
\author{Jiabin Qian, Manli Song}
\date{Source: arXiv:2607.23578v1 (CC BY 4.0)}
\begin{document}
\maketitle

\noindent\emph{Source and licence.} Decay estimates for a class of dispersive equations with partial inverse-square potentials. Version arXiv:2607.23578v1 (CC BY 4.0), distributed by arXiv under the Creative Commons Attribution 4.0 International licence, http://creativecommons.org/licenses/by/4.0/, as recorded in the arXiv licence metadata for this exact version and shown on the abstract page https://arxiv.org/abs/2607.23578v1. Full licence notice: \texttt{licenses/arxiv-2607.23578-CC-BY-4.0.txt}.\par\smallskip

\noindent\textit{Editorial scope.} Counted unit: the article's lemma lemma (source0.tex lines 799--809). The article's complete original proof of it is delivered with it (source0.tex lines 810--944, one \texttt{proof} environment of 135 lines). No mathematics is written, repaired or re-derived: the statement and the proof are the article's own lines, in the article's own notation, and no retained line is substituted or re-flowed.  Retained uncounted as the source-local context the proof needs: lines 251--279; lines 484--490; lines 523--529.  Nothing was omitted from any retained span, so the package carries no omission record.  No retained line contains a citation, so the wrapper carries no bibliography and no bibliography entry.  No retained line contains a figure environment, an image-inclusion command or drawing code, so no image is delivered and the package declares no asset dependency.  Editorial formatting only, in addition: an AMS \texttt{amsart} preamble that restates the article's own declarations and macros used by the retained lines, namely the environments \texttt{theorem}, \texttt{lemma} on separate counters, together with the standard packages those lines need.  The retained environments print in this document as Theorem 1, Lemma 1 in order of appearance, whereas the article prints the same items with the numbering of its own full document.  The complete native source of the article as distributed in the arXiv e-print for this exact version is bundled unchanged as \texttt{sources/206038/source0.tex}.
\begin{theorem}\label{ResultTime}
	Assume $\phi:\mathbb{R}^+ \to \mathbb{R}$ is smooth. Let  $U_k(t)=e^{it\phi(\sqrt{\mathcal{L}_a})}\psi(2^{-k}\sqrt{\mathcal{L}_a}), \forall k\in\mathbb{Z}$ and $U^{low}(t)=e^{it\phi(\sqrt{\mathcal{L}_a})}\varphi(\sqrt{\mathcal{L}_a})=\sum\limits_{k\leq0} U_k(t)$.   Then the subsequent outcomes are valid.
 \begin{enumerate}
 \item[\textbf{(A)}]   For $k \geq 1$, assume that $\phi$ satisfies (C1), then
	\begin{equation*}\label{res3-2}
	\|U_k(t) f\|_{L^\infty(\mathbb{R}^{2+n})} \lesssim |t|^{-\theta}2^{k\left(n+2-m_1\theta\right)}\|f\|_{L^1(\mathbb{R}^{2+n})},\,0 \leq \theta \leq \frac{n+1}{2}.
	\end{equation*}
 In addition, if $\phi$ satisfies (C3), then
	\begin{equation*}\label{res3-3}
	\|U_k(t) f\|_{L^\infty(\mathbb{R}^{2+n})} \lesssim |t|^{-\frac{n+1+\theta}{2}}2^{k\left(n+2-\frac{m_1(n+1+\theta)}{2}-\frac{\theta(\alpha_1-m_1)}{2}\right)}\|f\|_{L^1(\mathbb{R}^{2+n})},\,   0 \leq \theta \leq 1.
	\end{equation*}
 \item[\textbf{(B)}] For $k\leq 0$, suppose $\phi$ satisfies (C2), then
	\begin{equation*}
	\|U_k(t) f\|_{L^\infty(\mathbb{R}^{2+n})} \lesssim |t|^{-\theta}2^{k\left(n+2-m_2\theta\right)}\|f\|_{L^1(\mathbb{R}^{2+n})},\,0 \leq \theta \leq \frac{n+1}{2}.
	\end{equation*}
In addition, if $\phi$ satisfies (C4), then
	\begin{equation*}
	\|U_k(t) f\|_{L^\infty(\mathbb{R}^{2+n})} \lesssim |t|^{-\frac{n+1+\theta}{2}}2^{k\left(n+2-\frac{m_2(n+1+\theta)}{2}-\frac{\theta(\alpha_2-m_2)}{2}\right)}\|f\|_{L^1(\mathbb{R}^{2+n})},\,   0 \leq \theta \leq 1.
	\end{equation*}
\item[\textbf{(C)}] If $\phi$ satisfies (C2), then
	\begin{equation}\label{res-sum1}
 \|U^{low}(t) f\|_{L^\infty(\mathbb{R}^{2+n})} \lesssim (1+|t|)^{-\theta}\|f\|_{L^1(\mathbb{R}^{2+n})},\,\theta=\min\left(\frac{n+2}{m_2},\frac{n+1}{2}\right).
	\end{equation}
In addition, if $\phi$ satisfies (C4) and $\alpha_2=m_2$, then
	\begin{equation}\label{res-sum2}
	\|U^{low}(t) f\|_{L^\infty(\mathbb{R}^{2+n})}  \lesssim (1+|t|)^{-\theta}\|f\|_{L^1(\mathbb{R}^{2+n})},\,\theta=\min\left(\frac{n+2}{m_2},\frac{n+2}{2}\right).
	\end{equation}
	 \end{enumerate}
\end{theorem}

\begin{equation}\label{spectral}
\begin{aligned}
dE_{\sqrt{\mathcal{L}_a}}(\lambda; x, y; x', y') &= \frac{\lambda^{n+1}}{\pi^{n+1}} \sum_{\pm} a_{\pm}(\lambda|(x, y) - (x', y')|) e^{\pm i\lambda|(x, y) - (x', y')|}  \\
 &+ \frac{\lambda^{n+1}}{\pi^{n+1}} \int_{|\theta - \theta'|}^{\pi} \sum_{\pm} a_{\pm}(\lambda|\vec{n}_s^1|) e^{\pm i\lambda|\vec{n}_s^1|} \times A_a(s; \theta, \theta') ds  \\
&- \frac{\lambda^{n+1}}{\pi^{n+1}} \int_0^\infty \sum_{\pm} a_{\pm}(\lambda|\vec{n}_s^2|) e^{\pm i\lambda|\vec{n}_s^2|} \times B_a(s, \theta, \theta') ds, 
\end{aligned}
\end{equation}

\begin{lemma}\label{Sum}
Fix $\beta>0$. There exists $C_\beta>0$ such that for any $A>0$, we have
\begin{align*}
\sum_{j\in\mathbb{Z}, 2^j\leq A}2^{j\beta}&\leq C_\beta A^\beta,\\
\sum_{j\in\mathbb{Z}, 2^j>A}2^{-j\beta}&\leq C_\beta A^{-\beta}.
\end{align*}
\end{lemma}

\begin{lemma}\label{key lemma}
Let $\nu>0$, $n\geq 1$ and $2<q<\infty$. Then we have
\begin{equation}
\|e^{it\mathcal{L}_a^\nu} \mathcal{L}_a^{\frac{2\nu}{q}-\frac{n+2}{2}}f\|_{L_{x,y}^\infty(\mathbb{R}^{2+n})}\lesssim |t|^{-\frac{2}{q}}\|f\|_{L_{x,y}^1(\mathbb{R}^{2+n})},
\end{equation}
if either of the following conditions holds
\begin{itemize}
\item $\nu=\frac{1}{2}$ and $\frac{2}{q}<\frac{n+1}{2}$.
\item $\nu\not=\frac{1}{2}$ and $\frac{2}{q}\leq\frac{n+2}{2}$.
\end{itemize}
\end{lemma}

\begin{proof}
By Theorem \ref{ResultTime}, we have for $0\leq \theta\leq \frac{n+2-n_\nu}{2}$,
\begin{equation*}
\|e^{it\mathcal{L}_a^\nu} \psi(2^{-k}\sqrt{\mathcal{L}_a})f\|_{L_{x,y}^\infty(\mathbb{R}^{2+n})}\lesssim |t|^{-\theta}2^{k(n+2-2\nu\theta)}\|f\|_{L_{x,y}^1(\mathbb{R}^{2+n})},
\end{equation*}
and then
\begin{equation}\label{decay-local}
\|e^{it\mathcal{L}_a^\nu} \psi(2^{-k}\sqrt{\mathcal{L}_a})\mathcal{L}_a^{\frac{2\nu}{q}-\frac{n+2}{2}}f\|_{L_{x,y}^\infty(\mathbb{R}^{2+n})}\lesssim |t|^{-\theta}2^{2k\nu(\frac{2}{q}-\theta)}\|f\|_{L_{x,y}^1(\mathbb{R}^{2+n})}.
\end{equation}
If $\nu=\frac{1}{2}$ and $\frac{2}{q}<\frac{n+1}{2}$, or if $\nu\not=\frac{1}{2}$ and $\frac{2}{q}<\frac{n+2}{2}$, i.e., $\frac{2}{q}<\frac{n+2-n_\nu}{2}$, we get
\begin{align*}
&\quad \|e^{it\mathcal{L}_a^\nu} \mathcal{L}_a^{\frac{2\nu}{q}-\frac{n+2}{2}}f\|_{L_{x,y}^\infty(\mathbb{R}^{2+n})}\\
&\leq \sum_{k\in\mathbb{Z}} \|e^{it\mathcal{L}_a^\nu} \psi(2^{-k}\sqrt{\mathcal{L}_a})\mathcal{L}_a^{\frac{2\nu}{q}-\frac{n+2}{2}}f\|_{L_{x,y}^\infty(\mathbb{R}^{2+n})}\\
&\lesssim \left(\sum_{2^k\leq |t|^{-\frac{1}{2\nu}}}2^{\frac{4k\nu}{q}}+\sum_{2^k\geq|t|^{-\frac{1}{2\nu}}}|t|^{-\frac{n+2-n_\nu}{2}}2^{2k\nu(\frac{2}{q}-\frac{n+2-n_\nu}{2})}\right)\|f\|_{L_{x,y}^1(\mathbb{R}^{2+n})}\\
&\lesssim |t|^{-\frac{2}{q}}\|f\|_{L_{x,y}^1(\mathbb{R}^{2+n})},
\end{align*}
where the second inequality comes from \eqref{decay-local} and the last inequality is obtained by applying Lemma \ref{Sum}. 

Therefore, it remains to prove the case: $\nu\not=\frac{1}{2}$ and $q=\frac{4}{n+2}$. Without loss of generality, we assume $t>0$.
We rewrite
\begin{equation*}
e^{it\mathcal{L}_a^\nu} \mathcal{L}_a^{\frac{2\nu}{q}-\frac{n+2}{2}}f(x,y)=\sum_{k\in\mathbb{Z}}e^{it\mathcal{L}_a^\nu} \psi(2^{-k}t^\frac{1}{2\nu}\sqrt{\mathcal{L}_a})\mathcal{L}_a^{\frac{2\nu}{q}-\frac{n+2}{2}}f(x,y),
\end{equation*}
where 
\begin{align*}
&\quad e^{it\mathcal{L}_a^\nu} \psi(2^{-k}t^\frac{1}{2\nu}\sqrt{\mathcal{L}_a})\mathcal{L}_a^{\frac{2\nu}{q}-\frac{n+2}{2}}f(x,y)\\
&=\int_{\mathbb{R}^{2+n}}\int_0^\infty e^{it\lambda^{2\nu}} \psi(2^{-k}t^\frac{1}{2\nu}\lambda)\lambda^{\frac{4\nu}{q}-(n+2)}dE_{\sqrt{\mathcal{L}_a}}(\lambda; x,y; x',y')f(x',y')dx'dy',
\end{align*}
with the kernel
\begin{equation*}
J_k(x,y;x',y')=\int_0^\infty e^{it\lambda^{2\nu}} \psi(2^{-k}t^\frac{1}{2\nu}\lambda)\lambda^{\frac{4\nu}{q}-(n+2)}dE_{\sqrt{\mathcal{L}_a}}(\lambda; x,y; x',y').
\end{equation*}
By the triangle inequality, it suffices to prove
\begin{equation}\label{key kernel}
\sum_{k\in\mathbb{Z}}|J_k(x,y;x',y')|\lesssim |t|^{-\frac{2}{q}}, \quad\forall (x,y),(x',y')\in\mathbb{R}^{2+n}.
\end{equation}
By \eqref{spectral}, let us divide the kernel $J_k$ into three parts by
\begin{equation*}
J_k(x,y;x',y')=\frac{1}{\pi^{n+1}}\sum_{\pm}\left(J_k^{\pm,1}(x-x',y-y')+J_k^{\pm,2}(x,y;x',y')+J_k^{\pm,3}(x,y;x',y')\right), 
\end{equation*}
where
\begin{equation*}
J_k^{\pm,1}(x,y)=\int_0^\infty  e^{it\lambda^{2\nu}} \psi(2^{-k}t^\frac{1}{2\nu}\lambda)\lambda^{\frac{4\nu}{q}-1}a_{\pm}(\lambda|(x,y)|)e^{\pm i\lambda|(x,y)|}d\lambda,
\end{equation*}
\begin{equation*}
J_k^{\pm,2}(x,y;x',y')=\int_{|\theta-\theta'|}^\pi \int_0^\infty   e^{it\lambda^{2\nu}} \psi(2^{-k}t^\frac{1}{2\nu}\lambda)\lambda^{\frac{4\nu}{q}-1}a_{\pm}(\lambda|\overset{\rightarrow}{n}_s^1|)e^{\pm i\lambda|\overset{\rightarrow}{n}_s^1|}d\lambda A_a(s;\theta,\theta')ds,
\end{equation*}
and
\begin{equation*}
J_k^{\pm,3}(x,y;x',y')=\int_0^\infty \int_0^\infty   e^{it\lambda^{2\nu}} \psi(2^{-k}t^\frac{1}{2\nu}\lambda)\lambda^{\frac{4\nu}{q}-1}a_{\pm}(\lambda|\overset{\rightarrow}{n}_s^2|)e^{\pm i\lambda|\overset{\rightarrow}{n}_s^2|}d\lambda B_a(s;\theta,\theta')ds.
\end{equation*}
Our aim \eqref{key kernel} becomes to prove the following estimates
\begin{align}
\sum_{k\in\mathbb{Z}}|J_k^{\pm,1}(x,y)|&\lesssim |t|^{-\frac{2}{q}},\label{1}\\
\sum_{k\in\mathbb{Z}}|J_k^{\pm,2}(x,y;x',y')|&\lesssim |t|^{-\frac{2}{q}},\label{2}\\
\sum_{k\in\mathbb{Z}}|J_k^{\pm,3}(x,y;x',y')|&\lesssim |t|^{-\frac{2}{q}},\label{3}
\end{align}
for any $(x,y),(x',y')\in\mathbb{R}^{2+n}$. We just prove \eqref{1}. In fact, replacing $|(x,y)|$ by $|\overset{\rightarrow}{n}_s^\ell|, \ell=2,3$ and arguing similarly as \eqref{1}, we can prove \eqref{2} and \eqref{3}. We rewrite $J_k^{\pm,1}$ by
\begin{equation*}
J_k^{\pm,1}(x,y)=|t|^{-\frac{2}{q}}\mathcal{J}_k^{\pm,1}(t^{-\frac{1}{2\nu}}x,t^{-\frac{1}{2\nu}}y)
\end{equation*}
where
\begin{align*}
\mathcal{J}_k^{\pm,1}(x,y)&=2^\frac{4k\nu}{q}\int_0^\infty  e^{i2^{2k\nu}\lambda^{2\nu}} \psi(\lambda)\lambda^{\frac{4\nu}{q}-1}a_{\pm}(2^k\lambda|(x,y)|)e^{\pm i2^k\lambda|(x,y)|}d\lambda\\
&=2^\frac{4k\nu}{q}\int_0^\infty  e^{i2^{2k\nu}\lambda^{2\nu}} \widetilde{\psi}(\lambda)a_{\pm}(2^k\lambda|(x,y)|)e^{\pm i2^k\lambda|(x,y)|}d\lambda,
\end{align*}
with $\widetilde{\psi}(\lambda)=\psi(\lambda)\lambda^{\frac{4\nu}{q}-1}$ such that $\supp \widetilde{\psi}\subseteq [1/2,1]$.

Therefore, the estimate \eqref{1} is equivalent to
\begin{equation}\label{J}
\sum_{k\in\mathbb{Z}}|\mathcal{J}_k^{\pm,1}(x,y)|\lesssim 1,\quad \forall (x,y)\in\mathbb{R}^{2+n}.
\end{equation}
As the proof of part \textbf{(A)} in Theorem \ref{ResultTime}, we shall discuss the above statement in two cases.

\textbf{\underline{Case 1.}} $2^k|(x,y)|\leq 1$. For any $Q\in\mathbb{N}$, using integration by part $Q$-times, we have
\begin{align*}
|\mathcal{J}_k^{\pm,1}(x,y)|&=\left|2^\frac{4k\nu}{q}\int_0^\infty  \left[\left(\frac{\partial_\lambda}{i2\nu 2^{2k\nu}\lambda^{2\nu-1}}\right)^Q e^{i2^{2k\nu}\lambda^{2\nu}}\right] \widetilde{\psi}(\lambda)a_{\pm}(2^k\lambda|(x,y)|)e^{\pm i2^k\lambda|(x,y)|}d\lambda\right|\\
&\lesssim 2^\frac{4k\nu}{q}2^{-2kQ\nu}.
\end{align*}
Hence, choosing $Q$ sufficiently large and applying Lemma \ref{Sum}, it yields
\begin{equation}\label{Case 1}
\begin{aligned}
\sum_{2^k|(x,y)|\leq 1}|\mathcal{J}_k^{\pm,1}(x,y)|&=\sum_{\substack{k\geq 1\\ 2^k|(x,y)|\leq 1}}|\mathcal{J}_k^{\pm,1}(x,y)|+\sum_{\substack{k\leq 0\\ 2^k|(x,y)|\leq 1}}|\mathcal{J}_k^{\pm,1}(x,y)|\\
&\lesssim \sum_{k\geq1}2^\frac{4k\nu}{q}2^{-2kQ\nu}+\sum_{k\leq 0}2^\frac{4k\nu}{q}\\
&\lesssim 1.
\end{aligned}
\end{equation}
\textbf{\underline{Case 2.}} $2^k|(x,y)|>1$. We rewrite $\mathcal{J}_k^{\pm,1}(x,y)$ by
\begin{align*}
\mathcal{J}_k^{\pm,1}(x,y)
&=2^\frac{4k\nu}{q}\int_0^\infty  e^{i\Phi^\pm_k(\lambda;x,y)}\widetilde{\psi}(\lambda)a_{\pm}(2^k\lambda|(x,y)|)d\lambda,
\end{align*} 
with $\Phi^\pm_k(\lambda;x,y)=2^{2k\nu}\lambda^{2\nu} \pm 2^k\lambda|(x,y)|$. We continue our discussion by considering the following three sub-cases:
\begin{itemize}
\item $2^{k(2\nu-1)}\sim |(x,y)|>2^{-k}$. In this case, $2^k|(x,y)|\sim 2^{2k\nu}$ and $|\partial_\lambda^2\Phi^\pm_k(\lambda;x,y)|\sim 2^{2k\nu}$. By the Van der Corput lemma, it indicates that
\begin{equation}
|\mathcal{J}_k^{\pm,1}(x,y)|\lesssim 2^\frac{4k\nu}{q} \left(2^{2k\nu}\right)^{-\frac{1}{2}}\left(2^k|(x,y)|\right)^{-\frac{n+1}{2}}\sim 2^{k\nu(\frac{4}{q}-(n+2))}=1.
\end{equation}
Hence, we have
\begin{equation}\label{Subcase 1}
\sum_{\substack{2^{k(2\nu-1)}\sim |(x,y)|\\ 2^k|(x,y)|>1}}|\mathcal{J}_k^{\pm,1}(x,y)|\lesssim 1.
\end{equation}
\item $|(x,y)|\gg 2^{k(2\nu-1)}$. In this case, $|\partial_\lambda\Phi^\pm_k(\lambda;x,y)|\gtrsim 2^k|(x,y)|\gg 2^{2k\nu}$. For any $Q\in\mathbb{N}$, using integration by part $Q$-times, we have for any $Q\in\mathbb{N}$,
\begin{equation}
\begin{aligned}
|\mathcal{J}_k^{\pm,1}(x,y)|&=\left|2^\frac{4k\nu}{q}\int_0^\infty  \left[\left(\frac{\partial_\lambda}{i\partial_\lambda\Phi^\pm_k(\lambda;x,y)}\right)^Q e^{i\Phi^\pm_k(\lambda;x,y)}\right] \widetilde{\psi}(\lambda)a_{\pm}(2^k\lambda|(x,y)|)d\lambda\right|\\
&\lesssim 2^\frac{4k\nu}{q}\left(2^k|(x,y)|\right)^{-Q-\frac{n+1}{2}}\lesssim 2^\frac{4k\nu}{q}2^{-2k\nu(Q+\frac{n+1}{2})}.
\end{aligned}
\end{equation}
Hence, argued as \eqref{Case 1}, we have
\begin{equation}\label{Subcase 2}
\sum_{\substack{|(x,y)|\gg 2^{k(2\nu-1)}\\ 2^k|(x,y)|>1}}|\mathcal{J}_k^{\pm,1}(x,y)|\lesssim 1.
\end{equation}
\item $|(x,y)|\ll 2^{k(2\nu-1)}$. In this case, $1<2^k|(x,y)|\ll 2^{2k\nu}$ yields $k\geq 0$. Besides, $|\partial_\lambda\Phi^\pm_k(\lambda;x,y)|\gtrsim 2^{2k\nu}$. For any $Q\in\mathbb{N}$, using integration by part $Q$-times, we have for any $Q\in\mathbb{N}$
\begin{equation}
\begin{aligned}
|\mathcal{J}_k^{\pm,1}(x,y)|&=\left|2^\frac{4k\nu}{q}\int_0^\infty  \left[\left(\frac{\partial_\lambda}{i\partial_\lambda\Phi^\pm_k(\lambda;x,y)}\right)^Q e^{i\Phi^\pm_k(\lambda;x,y)}\right] \widetilde{\psi}(\lambda)a_{\pm}(2^k\lambda|(x,y)|)d\lambda\right|\\
&\lesssim 2^\frac{4k\nu}{q}\left(2^{2k\nu}\right)^{-Q}\left(2^k|(x,y)|\right)^{-\frac{n+1}{2}}\lesssim 2^\frac{4k\nu}{q}\left(2^{2k\nu}\right)^{-Q}.
\end{aligned}
\end{equation}
Hence, argued as \eqref{Case 1}, we have
\begin{equation}\label{Subcase 3}
\sum_{\substack{|(x,y)|\ll 2^{k(2\nu-1)}\\ 2^k|(x,y)|>1}}|\mathcal{J}_k^{\pm,1}(x,y)|\lesssim 1.
\end{equation}
\end{itemize}
Therefore, by \eqref{Subcase 1}, \eqref{Subcase 2} and \eqref{Subcase 3}, we obtain
\begin{equation}\label{Case 2}
\begin{aligned}
&\quad\sum_{2^k|(x,y)|>1}|\mathcal{J}_k^{\pm,1}(x,y)|\\
&=\sum_{\substack{2^{k(2\nu-1)}\sim |(x,y)|\\ 2^k|(x,y)|>1}}|\mathcal{J}_k^{\pm,1}(x,y)|+\sum_{\substack{|(x,y)|\gg 2^{k(2\nu-1)}\\ 2^k|(x,y)|>1}}|\mathcal{J}_k^{\pm,1}(x,y)|+\sum_{\substack{|(x,y)|\ll 2^{k(2\nu-1)}\\ 2^k|(x,y)|>1}}|\mathcal{J}_k^{\pm,1}(x,y)|\\
&\lesssim 1.
\end{aligned}
\end{equation}
The estimate \eqref{J} comes immediately from \eqref{Case 1} and \eqref{Case 2}, which completes our proof.\\
\end{proof}

\end{document}
